% ============================================================================
%  ReqGuard AI -- Project Documentation (6-page condensed version)
%  Delulu 2 Deploy Hackathon -- Problem Statement G14 (Generative AI)
%
%  Compile:  pdflatex ReqGuard_AI_Documentation_6pages.tex   (run twice)
%
%  ---------------------------------------------------------------------------
%  PAGE-FIT KNOBS -- if it lands at 5 or 7 pages, adjust in this order:
%    1. \documentclass[10pt] -> [11pt] to grow, or margin=0.85in -> 0.95in
%    2. The \compactspacing block below (\parskip / \itemsep)
%    3. Drop the Listing 1 code block (Sec. 6.2) to lose ~1/3 page
%    4. Scale the figures: scale=0.92 in the two \begin{tikzpicture}[...]
%  ---------------------------------------------------------------------------
%  EDIT BEFORE SUBMITTING: \TeamName, \TeamMembers, \CollegeName
% ============================================================================
\documentclass[10pt,a4paper]{article}

\usepackage[utf8]{inputenc}
\usepackage[T1]{fontenc}
\usepackage{lmodern}
\usepackage{microtype}

\usepackage[margin=0.85in,top=0.8in,bottom=0.8in]{geometry}
\usepackage{titlesec}
\usepackage{enumitem}
\usepackage{xcolor}
\usepackage{booktabs}
\usepackage{array}
\usepackage{listings}
\usepackage{caption}
\usepackage{fancyhdr}
\usepackage{tikz}
\usetikzlibrary{arrows.meta, positioning, fit, backgrounds}
\usepackage{hyperref}

\definecolor{primary}{RGB}{79,70,229}
\definecolor{accent}{RGB}{13,148,136}
\definecolor{warn}{RGB}{180,83,9}
\definecolor{codebg}{RGB}{247,247,250}
\definecolor{codeframe}{RGB}{205,205,215}

\hypersetup{colorlinks=true, linkcolor=primary, urlcolor=accent,
  pdftitle={ReqGuard AI -- Project Documentation},
  pdfsubject={Delulu 2 Deploy Hackathon -- Problem Statement G14}}

% ---------- compact spacing ----------
\setlength{\parskip}{3pt}
\setlength{\parindent}{0pt}
\setlist{nosep, leftmargin=1.25em, topsep=2pt, itemsep=1.5pt}
\captionsetup{font=small, skip=4pt}
\titlespacing*{\section}{0pt}{8pt}{3pt}
\titlespacing*{\subsection}{0pt}{6pt}{2pt}
\titleformat{\section}{\large\bfseries\color{primary}}{\thesection}{0.5em}{}
\titleformat{\subsection}{\normalsize\bfseries\color{primary!80!black}}{\thesubsection}{0.5em}{}

\pagestyle{fancy}
\fancyhf{}
\fancyhead[L]{\footnotesize ReqGuard AI --- Requirements Engineering Agent (G14)}
\fancyhead[R]{\footnotesize \thepage}
\renewcommand{\headrulewidth}{0.4pt}

\lstdefinestyle{code}{backgroundcolor=\color{codebg}, basicstyle=\ttfamily\scriptsize,
  breaklines=true, frame=single, rulecolor=\color{codeframe}, framesep=4pt,
  columns=fullflexible, keepspaces=true, tabsize=2, showstringspaces=false,
  aboveskip=5pt, belowskip=4pt}
\lstset{style=code}

% ---------- EDIT THESE ----------
\newcommand{\TeamName}{Byteforce}
\newcommand{\TeamMembers}{Member 1, Member 2, Member 3, Member 4}
\newcommand{\CollegeName}{[Your College / Institution]}

\begin{document}

% ================= TITLE BLOCK (compact, no separate title page) =============
\begin{center}
  {\LARGE\bfseries\color{primary} ReqGuard AI}\\[2pt]
  {\large\itshape An AI Requirements Engineering Agent}\\[4pt]
  {\footnotesize Delulu 2 Deploy Hackathon \quad$\vert$\quad Problem Statement
   \textbf{G14} (Generative AI)}\\[4pt]
  {\footnotesize \textbf{Team:} \TeamName \quad$\vert$\quad \TeamMembers
   \quad$\vert$\quad \CollegeName \quad$\vert$\quad \today}
\end{center}
\vspace{2pt}
\hrule
\vspace{6pt}

\begin{center}
\begin{minipage}{0.93\textwidth}
{\small\textbf{Abstract.} Requirements documents are where software projects go
wrong first: a requirement that says ``the application shall respond quickly''
cannot be tested, and two requirements that each claim exclusive control of
login cannot both be built. ReqGuard AI reads a Software Requirements
Specification (PDF, DOCX or TXT), extracts the individual requirements, and
reports the ambiguities, contradictions, duplicates, missing details and hidden
dependencies inside it. Every finding quotes evidence from the document and
traces back to the requirement it came from. The system is a complete
end-to-end application --- a modular pipeline of independently testable stages,
not a single prompt --- verified by 73 unit tests, 49 live API checks and 28
browser checks, with evaluation figures measured from a real run rather than
asserted.}
\end{minipage}
\end{center}
\vspace{4pt}

% ============================================================================
\section{Problem and Objectives}

A requirement that cannot be tested, one that is silently repeated in two
sections, or two that cannot both be satisfied will each become a defect, a
duplicated implementation, or a late redesign. None of this is visible until
somebody reads the entire specification carefully, and most teams never do.
Problem statement G14 asks for a generative-AI agent that performs this review
automatically, before implementation begins.

\textbf{Objectives.} (i)~Extract individual requirements from PDF/DOCX/TXT while
preserving exact original wording; (ii)~detect ambiguous, untestable
requirements and quote the offending term; (iii)~detect semantic duplicates;
(iv)~detect contradictory pairs with the conflicting clause from each side;
(v)~detect missing information an implementer would have to ask about;
(vi)~build a requirement dependency graph; (vii)~produce a traceability matrix
linking every finding to its source; (viii)~let a reviewer accept, edit or
reject each suggested refinement and export a clean requirement set without
ever discarding the original text; and (ix)~deliver all of this as a real,
deployable application --- no mockups, no hard-coded findings, no fabricated
statistics.

\begin{table}[h!]
\centering\small
\begin{tabular}{@{}p{3.1cm} p{12.3cm}@{}}
\toprule
\textbf{Output} & \textbf{Description} \\
\midrule
Clean requirement set & Every requirement with accepted refinements applied and confirmed duplicates marked; original wording always retained. \\
Ambiguity report & Requirements an engineer and a tester could read differently, with the offending term quoted. \\
Contradiction report & Requirement pairs that cannot both be satisfied, side by side with the conflicting clause from each. \\
Duplicate report & Pairs confirmed to state the same requirement, with the embedding similarity that surfaced them. \\
Dependency graph & An interactive graph of which requirements must be built before others. \\
Traceability matrix & Every requirement, its source section/page, and every finding raised against it. \\
\bottomrule
\end{tabular}
\caption{What the system produces}
\end{table}

% ============================================================================
\section{System Architecture}

The frontend never talks to the database and never holds an API key. Every
credential lives server-side; the browser only calls the FastAPI service over
HTTPS.

\begin{figure}[h!]
\centering
\begin{tikzpicture}[scale=1, transform shape,
  node distance=0.75cm and 1.4cm,
  box/.style={draw, rounded corners, thick, minimum height=0.8cm, minimum width=2.9cm, align=center, font=\scriptsize},
  ui/.style={box, fill=primary!10, draw=primary},
  be/.style={box, fill=accent!12, draw=accent},
  svc/.style={box, fill=primary!5, draw=primary!60, minimum width=3.5cm},
  data/.style={box, fill=warn!8, draw=warn!70, minimum width=3.5cm},
  arr/.style={-{Latex[length=2mm]}, thick}]

  \node[ui] (front) {Browser\\React 19 + Vite + TS};
  \node[be, right=2.0cm of front] (api) {FastAPI Backend\\Python 3.12};
  \node[svc, above right=0.35cm and 1.5cm of api] (emb) {EmbeddingService\\BAAI/bge-small-en-v1.5 (384d)};
  \node[svc, right=1.5cm of api] (ai) {AIService\\Groq LLM};
  \node[data, below right=0.35cm and 1.5cm of api] (db) {InsForge PostgreSQL\\+ pgvector (HNSW)};
  \node[data, below=0.75cm of ai] (store) {InsForge Storage};

  \draw[arr] (front) -- node[above, font=\tiny]{HTTPS / REST} (api);
  \draw[arr] (api) -- (emb);
  \draw[arr] (api) -- (ai);
  \draw[arr] (api) -- (db);
  \draw[arr] (api) -- (store);
  \begin{scope}[on background layer]
    \node[draw=black!30, dashed, rounded corners, inner sep=6pt, fit=(api)(emb)(ai)] {};
  \end{scope}
\end{tikzpicture}
\caption{Architecture. All secrets, model access and database access are server-side only.}
\end{figure}

% ============================================================================
\section{The Analysis Pipeline}

Each stage below is a separate service with its own prompt and its own tests.
There is deliberately \textbf{no single prompt that ``does the analysis''}, so
one stage can fail without corrupting the others, and each can be tested
independently.

\begin{figure}[h!]
\centering
\begin{tikzpicture}[scale=0.97, transform shape,
  node distance=0.5cm and 0.75cm,
  st/.style={draw, rounded corners, thick, minimum height=0.72cm, minimum width=2.15cm, align=center, font=\tiny, fill=primary!8, draw=primary!70},
  det/.style={draw, rounded corners, thick, align=left, font=\tiny, fill=accent!10, draw=accent!70, inner sep=5pt},
  out/.style={st, fill=warn!10, draw=warn!70},
  arr/.style={-{Latex[length=1.8mm]}, thick}]

  \node[st] (doc) {Document\\PDF / DOCX / TXT};
  \node[st, right=of doc] (parse) {Parser\\page + section};
  \node[st, right=of parse] (ext) {Extraction\\original wording};
  \node[st, right=of ext] (norm) {Normalisation\\+ embeddings};
  \node[st, right=of norm] (cand) {Candidate pairs\\cosine similarity};

  \node[det, below=0.85cm of norm, text width=7.4cm] (det)
    {\textbf{Five independent detection services} (each: own prompt, own schema, own tests)\\
     Duplicate verification \;$\vert$\; Contradiction verification \;$\vert$\; Ambiguity\\
     Missing information \;$\vert$\; Dependency analysis};

  \node[out, below=0.75cm of det, minimum width=2.6cm] (ref) {Refinement generation};
  \node[out, left=0.9cm of ref, minimum width=2.6cm] (trace) {Traceability +\\clean set};
  \node[out, right=0.9cm of ref, minimum width=2.6cm] (dash) {Database\\$\rightarrow$ Dashboard};

  \draw[arr] (doc) -- (parse);
  \draw[arr] (parse) -- (ext);
  \draw[arr] (ext) -- (norm);
  \draw[arr] (norm) -- (cand);
  \draw[arr] (cand) |- (det.east);
  \draw[arr] (ext) |- (det.west);
  \draw[arr] (det) -- (ref);
  \draw[arr] (ref) -- (trace);
  \draw[arr] (ref) -- (dash);
\end{tikzpicture}
\caption{The pipeline. Pair-wise stages are fed by similarity-ranked candidates; per-requirement stages read the extracted text directly.}
\end{figure}

\textbf{Why two-stage detection matters.} Comparing every requirement against
every other is quadratic: 100 requirements is 4{,}950 pairs, and sending all of
them to a language model is slow and expensive. Instead, embeddings rank pairs
by meaning and only those above a configurable similarity threshold reach the
model. On the sample specification this cuts 1{,}128 possible pairs to 16
duplicate candidates and 265 contradiction candidates. Similarity never decides
anything by itself --- it only chooses what is worth reasoning about; the model
makes the call and must quote evidence for it.

\begin{table}[h!]
\centering\small
\begin{tabular}{@{}p{2.4cm} p{13.0cm}@{}}
\toprule
\textbf{Layer} & \textbf{Technologies} \\
\midrule
Frontend & React 19, TypeScript, Vite, Tailwind CSS v4, React Flow (\texttt{@xyflow/react}) for the dependency graph, React Router, Lucide \\
Backend & Python 3.12+, FastAPI, Uvicorn, Pydantic v2, pydantic-settings \\
AI & Groq (\texttt{openai/gpt-oss-20b}); an OpenRouter-compatible client sits behind the same interface, selected by \texttt{LLM\_PROVIDER} \\
Embeddings & \texttt{BAAI/bge-small-en-v1.5} (sentence-transformers), 384 dimensions, cosine similarity \\
Database & InsForge PostgreSQL with \texttt{pgvector} and an HNSW cosine index \\
Parsing & pypdf, python-docx, plain text \\
\bottomrule
\end{tabular}
\caption{Technology stack}
\end{table}

% ============================================================================
\section{Implementation}

\subsection{One choke point for every model call}

Every language-model call in the codebase passes through a single class,
\texttt{AIService}. Each call requests a JSON object validated against a
Pydantic schema; invalid JSON triggers exactly one retry with a correction
prompt; if that also fails, the stage raises and the run is marked
\texttt{failed} with the real reason recorded. The system never substitutes a
fabricated result to fill a gap.

\subsection{Case study: the RequirementResolver bug}

The first full end-to-end run reported \textbf{zero} duplicates and
\textbf{zero} contradictions. The pipeline was not broken --- the model was
returning the full requirement \emph{sentence} where the prompt asked for a
short code (\texttt{R001}), and a strict dictionary lookup silently discarded
every verdict that did not match a code exactly. A dedicated resolver now
accepts four forms of reference:

\begin{lstlisting}[language=Python, caption={\texttt{app/utils/resolve.py} --- defensive entity-reference resolution}]
def resolve(self, reference: str | None) -> HasCodeAndText | None:
    direct = self._by_code.get(raw.upper())        # 1. exact code, the documented contract
    if direct is not None:
        return direct
    match = _CODE_PATTERN.search(raw)              # 2. code embedded in text ("Requirement R001")
    if match:
        found = self._by_code.get(match.group().upper())
        if found is not None:
            return found
    exact_text = self._by_text.get(_key(raw))      # 3. full sentence echoed back instead of a code
    if exact_text is not None:
        return exact_text
    if len(probe) >= 25:                           # 4. truncated or lightly reworded quote
        for text, record in self._by_text.items():
            if probe in text or text in probe:
                return record
    return None
\end{lstlisting}

After wiring this into the duplicate, contradiction, ambiguity, dependency and
refinement services, the same document went from 0 to 2 duplicates and from 0
to 40 contradictions. The findings had been there all along, only unread.

\subsection{Thresholds, schema and security}

\textbf{Thresholds are configuration, not constants,} because no single value is
right for every document. Defaults were chosen by measuring the similarity
distribution of the sample, not guessed: duplicate candidate cutoff
\texttt{0.82}, contradiction cutoff \texttt{0.70}, candidate cap \texttt{300},
and per-class confidence floors (contradiction \texttt{0.60}, partial conflict
\texttt{0.85}, duplicate \texttt{0.60}, dependency \texttt{0.40}).

\textbf{Schema.} Seven tables --- \texttt{documents}, \texttt{requirements},
\texttt{analysis\_runs}, \texttt{issues}, \texttt{issue\_relationships},
\texttt{dependencies}, \texttt{refinements} --- all with UUID primary keys, and
a \texttt{vector(384)} embedding column indexed with HNSW.

\textbf{Security.} Secrets exist only in backend environment variables:
\texttt{GROQ\_API\_KEY} and the InsForge admin key are never sent to the
browser, and no secret uses the \texttt{VITE\_} prefix, since anything with that
prefix is compiled into the client bundle. Row Level Security is enabled on all
seven tables with \emph{no} permissive policies and privileges revoked from the
\texttt{anon} and \texttt{authenticated} roles --- verified live: an anonymous
key returns HTTP 401 (\emph{permission denied for table documents}). Uploads are
validated for type, size and filename, and a content/extension mismatch is
rejected (415 for an unsupported type, 400 for an empty or mismatched file).
CORS uses an explicit origin list, never a wildcard. Errors are logged
server-side and returned as short, user-safe messages with no stack trace.
\texttt{.env} is git-ignored on both sides; a scan of tracked source found no
hardcoded credentials.

\begin{table}[h!]
\centering\small
\begin{tabular}{@{}p{1.2cm} p{7.4cm} p{6.6cm}@{}}
\toprule
\textbf{Method} & \textbf{Path} & \textbf{Purpose} \\
\midrule
GET & \texttt{/health}, \texttt{/api/system/status} & Liveness; real readiness of DB, embeddings, LLM \\
POST & \texttt{/api/documents/upload} & Upload an SRS (multipart \texttt{file}) \\
POST & \texttt{/api/documents/\{id\}/analyze} & Start the analysis run \\
GET & \texttt{/api/documents/\{id\}/status} & Stage, progress percentage, stats \\
GET & \texttt{/api/documents/\{id\}/requirements} & Extracted requirements \\
GET & \texttt{/api/documents/\{id\}/issues} & All findings; filters: type, severity, requirement, min.\ confidence \\
GET & \texttt{/api/documents/\{id\}/ambiguities} $\vert$ \texttt{contradictions} $\vert$ \texttt{duplicates} $\vert$ \texttt{missing-information} & The four per-type reports \\
GET & \texttt{/api/documents/\{id\}/dependencies} & Dependency graph nodes and edges \\
GET & \texttt{/api/documents/\{id\}/traceability} & Traceability matrix \\
GET & \texttt{/api/documents/\{id\}/clean-requirements} & Clean requirement set \\
GET & \texttt{/api/documents/\{id\}/export?format=json$\vert$markdown} & Export the clean set \\
POST & \texttt{/api/refinements/\{id\}/accept} $\vert$ \texttt{reject} $\vert$ \texttt{edit} & Act on a suggested rewrite \\
\bottomrule
\end{tabular}
\caption{REST API. Interactive OpenAPI documentation is served at \texttt{/docs}.}
\end{table}

% ============================================================================
\section{Testing and Verification}

The system was verified at three independent levels, each exercising the real
application rather than a mock of it.

\begin{table}[h!]
\centering\small
\begin{tabular}{@{}p{4.4cm} p{1.9cm} p{9.0cm}@{}}
\toprule
\textbf{Level} & \textbf{Result} & \textbf{Coverage} \\
\midrule
Unit tests (\texttt{pytest}) & 73 / 73 passed & Parsing, upload validation, embeddings, candidate selection, every detection stage, AI validation and retry, reference resolution, endpoints, traceability, refinement round-trip \\
Live demo flow (\texttt{verify\_demo\_flow.py}) & 49 / 49 passed & End-to-end walk against a running backend: health, rejection paths, a real analysis run, every report, filters, accept/edit/reject, export, 404s \\
Browser UI checks & 28 / 28 passed & Rendered pages against live data in an installed browser \\
\bottomrule
\end{tabular}
\caption{Verification summary}
\end{table}

Unit tests use a fake model client and an in-memory database, so they need no
network access and cost nothing; the other two levels run against the real
backend, database and model.

% ============================================================================
\section{Evaluation}

\texttt{samples/evaluation\_dataset.json} holds ground-truth labels for a
48-requirement sample specification containing deliberately seeded defects.
\texttt{backend/evaluate.py} scores a completed run against those labels.

\begin{table}[h!]
\centering\small
\begin{tabular}{@{}l c c c@{}}
\toprule
\textbf{Stage} & \textbf{Precision} & \textbf{Recall} & \textbf{F1} \\
\midrule
Requirement extraction & 1.00 & 1.00 & 1.00 \\
Ambiguity detection & 0.67 & 1.00 & 0.80 \\
Duplicate detection & 0.33 & 1.00 & 0.50 \\
Contradiction detection & 0.08 & 1.00 & 0.15 \\
Contradiction (confidence $\geq$ 0.9) & 0.20 & 1.00 & 0.33 \\
Missing information & 0.09 & 1.00 & 0.16 \\
\bottomrule
\end{tabular}
\caption{Measured on one run of the 48-requirement sample. Traceability: 76 of 76 issues carried a source requirement id (100\%).}
\end{table}

\textbf{Read these numbers carefully.} Recall is high --- every labelled problem
was found. Precision is low for two separate and identifiable reasons. First,
the ground-truth set is deliberately small (3 labelled contradictions, 2
labelled missing-information items), so many findings counted as false
positives are real problems that simply were never labelled. Second,
contradiction detection genuinely over-reports: a requirement containing an
exclusivity word such as ``only'' conflicts, in a technically correct sense,
with many others, producing a cascade of true-but-redundant pairs from one root
cause. These are measurements from one run on one small document, not a general
accuracy claim, and no figure shown anywhere in the application is hard-coded.

% ============================================================================
\section{Deployment}

Neither service assumes \texttt{localhost}: the frontend reads
\texttt{VITE\_API\_URL} at build time and the backend reads
\texttt{CORS\_ORIGINS} at runtime. \textbf{Backend --- Render:}
\texttt{backend/render.yaml} is a ready blueprint; \texttt{INSFORGE\_URL},
\texttt{INSFORGE\_API\_KEY}, \texttt{GROQ\_API\_KEY} and \texttt{CORS\_ORIGINS}
are set as dashboard secrets and health checks use \texttt{/health}. The
embedding model needs roughly 500\,MB of RAM, so the free instance tier is
insufficient and the blueprint specifies Starter. \textbf{Frontend --- Vercel:}
\texttt{vercel --prod} with \texttt{VITE\_API\_URL} pointed at the deployed
backend, whose \texttt{CORS\_ORIGINS} then includes the Vercel domain.

% ============================================================================
\section{Design Decisions and Limitations}

\textbf{Design decisions.}
\begin{itemize}
  \item \textbf{Original wording is never overwritten.} Extraction copies each requirement exactly as printed; normalisation produces a separate comparison string used only for embeddings; refinements live in their own table as proposals, and accepting one changes what the clean set renders, not the stored requirement.
  \item \textbf{Every finding must trace to a requirement.} Issues carry a requirement id and pair findings carry both. A finding that cannot be traced is not reported.
  \item \textbf{Models return codes, but not reliably} --- hence the resolver of Section~4.2, without which correct findings were silently discarded.
  \item \textbf{Failures are visible.} A stage that cannot parse a model response fails the run and records why; nothing is invented to fill the gap.
\end{itemize}

\textbf{Known limitations.}
\begin{itemize}
  \item \textbf{Contradiction precision is low.} One exclusivity clause generates many true-but-redundant pairs. The report offers a ``direct contradictions only'' filter and tunable confidence floors, but grouping cascading conflicts under a single root cause is not implemented.
  \item \textbf{Dependency recall varies between runs,} because requirements are analysed in overlapping windows and an edge between two distant requirements can be missed.
  \item \textbf{Analysis takes several minutes} for 48 requirements, dominated by sequential model calls; running the independent stages concurrently would cut this substantially.
  \item \textbf{The evaluation set is small} --- one document, a handful of labels. Treat the figures as a smoke test of behaviour, not a benchmark.
  \item \textbf{Uploaded bytes are held in memory} between upload and analysis, so a restart between the two steps requires re-uploading.
\end{itemize}

% ============================================================================
\section{Conclusion}

ReqGuard AI is a complete, working requirements-engineering agent: upload,
extraction, semantic similarity, duplicate detection, ambiguity detection,
contradiction detection, dependency analysis, traceability, refinement and a
real dashboard --- with no stage mocked and no statistic invented. Every claim
here is backed by a check that was actually run: 73 unit tests, 49 live API
checks and 28 browser checks, all passing, plus a measured evaluation against a
labelled sample. The clearest next steps are grouping cascading contradictions
under a single root cause to raise precision without losing recall,
parallelising the independent detection stages to cut analysis time, and
expanding the labelled evaluation set across more varied specifications.

\vspace{4pt}
{\footnotesize\textit{Built with Groq, InsForge (PostgreSQL + pgvector),
\texttt{BAAI/bge-small-en-v1.5} via sentence-transformers, FastAPI and React.}}

\end{document}
